\chapter{Introduction}
\label{ch:introduction}

This chapter introduces the problem the thesis addresses, motivates the choice to tackle it through architectural rather than incremental means, formalises the problem statement, declares the objectives, summarises the proposed system at a high level, and registers the scope within which the work is bounded.

\section{Background}
\label{sec:background}

\subsection{Capability and the persistent failure mode}
\label{sec:cap-failure}

Large language models trained on broad text corpora and instruction-tuned for assistant-style interaction now perform competitively on a wide range of natural language tasks, including factoid question answering, multi-step reasoning, summarisation, and program synthesis. Frontier systems such as GPT-4, Claude, and PaLM reach or exceed reported human accuracy on standardised academic benchmarks, and their accessibility through chat interfaces has made them part of everyday research, education, and professional workflows \cite{openai2023gpt4, anthropic2023claude, anil2023palm}. The same systems, however, generate confident but factually incorrect statements at deployment-relevant rates. The accepted name for this failure mode is hallucination, although the term covers heterogeneous causes ranging from retrieval limitations to confabulation under genuine uncertainty.

A practical consequence is that the same fluency that makes large language models useful also makes their hallucinations dangerous. A false answer phrased with the same confidence as a true one is harder to detect than a false answer that sounds uncertain. Users without independent verification, particularly in high-stakes domains, take the model's confident phrasing as a signal of reliability and then act on incorrect information. This thesis takes the position that hallucination is therefore not merely a benchmark-accuracy problem but a deployment-safety problem, and that addressing it requires architectural change rather than incremental decoding tricks.

\subsection{The hallucination problem at scale}
\label{sec:halluc-scale}

The empirical record on hallucination is consistent across benchmarks and domains. The TruthfulQA benchmark of Lin and colleagues isolates questions on which humans themselves are prone to hold confident misconceptions, and reports that the largest frontier model evaluated achieved only fifty-eight percent on the truthfulness axis and twenty-five percent on the combined truthful-and-informative axis, against a human baseline of ninety-four percent on both \cite{lin-etal-2022-truthfulqa}. The benchmark also documented an inverse-scaling relationship in which larger models were more, not less, prone to echo human falsehoods, indicating that scale alone does not solve the problem.

Long-form generation shows the same pattern at finer granularity. The atomic-fact decomposition protocol of Min and colleagues, evaluating chat-style models on biographical generation, found that only fifty-eight percent of decomposed atomic claims were supported by retrieved evidence, and the remaining claims ranged from unsupported speculation to outright fabrication \cite{min-etal-2023-factscore}. Domain-specific evaluations are still more concerning. Clinical-question studies of frontier models report hallucination rates between sixty-nine and seventy-seven percent of responses \cite{alkaissi2023artificial}, an order of magnitude above any threshold a responsible deployment would tolerate.

These results establish that hallucination is neither an artefact of small models nor a residual that scale will eliminate; it is, as the recent survey literature argues, a structural property of stateless generation under uncertainty rather than a defect to be trained away \cite{ji2023survey,huang-etal-2023-large}. Any system that aims to reduce it must change the structure rather than merely train harder.

\subsection{The stateless gap}
\label{sec:stateless-gap}

Three structural shortcomings of the current inference paradigm shape every other limitation a deployed model exhibits. First, standard inference is stateless: every query is processed independently, no representation of past successful answers is retained across sessions, and identical questions can produce inconsistent answers from one interaction to the next. This wastes computation on queries the system has effectively already solved and prevents the accumulation of domain-specific knowledge through use.

Second, modern neural networks are systematically overconfident. Calibration studies report that both classifiers and language models assign high probability to incorrect outputs at rates that exceed their actual accuracy by ten percentage points or more \cite{pmlr-v70-guo17a, kadavath2022language}. Single-signal correctness predictors achieve area-under-the-receiver-operating-curve scores between zero point five zero and zero point six zero, only marginally above random chance \cite{xiong2023can}. A user therefore cannot reliably tell which answers are likely to be wrong from confidence alone.

Third, deployed models are static. Once trained they neither learn from their mistakes nor adapt to the patterns of an institution that uses them, because naive online fine-tuning incurs catastrophic forgetting on the broad capabilities the base training instilled \cite{doi:10.1073/pnas.1611835114}. The implication is that verification effort spent today produces no compounding benefit tomorrow: a thousand answers verified by humans this morning will not change a single answer the model gives this afternoon.

These three shortcomings reinforce one another. Without memory the model cannot reuse verified knowledge. Without calibrated confidence the system cannot recognise when memory is missing. Without self-improvement the verifications that humans or external systems do produce decay into one-shot interventions. The remainder of this thesis develops a system that breaks the cycle along all three axes simultaneously. The next section argues why a combined intervention at the architectural level is the appropriate response.

\section{Rationale of the Study or Motivation}
\label{sec:rationale}

\subsection{The deployment gap}
\label{sec:deployment-gap}

The capability gap that motivates this thesis is the distance between what large language models do well and what reliable deployment requires of them. Fluency, breadth of topic coverage, and on-demand reasoning are already present in current frontier models. What is not present is the discipline of distinguishing facts that the model has correctly internalised from confabulations that emerge under genuine uncertainty, and the discipline of accumulating verified knowledge across interactions instead of regenerating answers from scratch every time \cite{farquhar2024semantic}. The gap is not aesthetic. In educational deployments, students who learn from confidently incorrect tutoring acquire misconceptions that persist after the source has been corrected. In healthcare information services, hallucinated drug interactions or treatment dosages produce decisions whose costs are measured in patient outcomes rather than in benchmark points. In enterprise knowledge workflows, fabricated citations and missing provenance erode the trust that the entire deployment rests on.

A second, less visible cost is computational. The stateless protocol of a typical large-model deployment regenerates every answer from first principles, including answers whose substance has already been verified in a previous interaction. As deployment scale grows, the cost of unnecessary regeneration accumulates against any computational budget. A system that can recognise a familiar query and serve a previously verified answer at near-zero cost converts the same compute envelope into more delivered value, and that conversion is itself part of the rationale.

\subsection{Why prior pipelines do not close the gap}
\label{sec:prior-pipelines}

Three families of prior systems have attempted to reduce hallucination, and each has a structural limitation that prevents it from closing the gap on its own. Retrieval-augmented generation grounds the model in retrieved evidence \cite{NEURIPS2020_6b493230}, which raises factual recall when the retrieval succeeds but is silent when retrieval misses or returns evidence the model summarises incorrectly. The output remains a single-pass generation whose quality is bounded by retrieval quality, and the system has no representation of which past answers were correct. Prompting interventions, including chain-of-thought reasoning and self-consistency over sampled chains \cite{wei2022chain, wang2023selfconsistency}, expose internal reasoning and exploit agreement across resampled answers, but they provide no persistence: each session begins with no memory of prior verified answers, and the same question asked tomorrow triggers the same reasoning effort as today. Post-generation verification \cite{manakul-etal-2023-selfcheckgpt, farquhar2024semantic} detects hallucinations after they occur and supports refusal or revision, but it does not feed back into model weights or memory, so the same hallucination patterns recur in the next session.

The implication is that hallucination reduction at deployment-relevant rates is unlikely to come from incremental improvement of any one family. Stronger retrievers, longer chains of thought, and sharper post-generation verifiers each yield diminishing returns when used in isolation. What is missing is a system in which verification feeds memory, memory feeds routing, routing feeds generation, and improvement compounds across cycles instead of evaporating at session boundaries. That system is an architectural intervention, and arguing for it is the work this thesis attempts.

\subsection{The architectural opportunity}
\label{sec:arch-opportunity}

Three capabilities, taken together, would close the gap that the existing three families leave open. The first is a verified episodic memory whose admission rule is finer than a single binary threshold. An answer that is plausible but has not earned the precision required for retrieval should not be discarded, because the same query may be answered more confidently after the model has been updated; it should instead enter a deferred buffer and be reconsidered at the next cycle boundary. A four-outcome decision (store, defer, abstain, discard) gives the system somewhere to put borderline outputs without contaminating the high-precision pool. The second capability is confidence-aware routing that reads pre-generation uncertainty as an independent safety signal alongside the combined memory similarity and stored quality, and that adds a per-query retrieval-utility decision on the residual queries that the combined-score dispatch sends to the retrieval-augmented fallback so that retrieval is consulted only on the subpopulation where it is expected to help rather than unconditionally. When the pre-generation uncertainty is high, the routing layer should force grounded retrieval before generation has started, rather than rely on post-generation review to catch an answer that should never have been produced. The third capability is a constrained self-improvement loop that fine-tunes on a verified high-quality pool under a regulariser pulling weights toward the base parameters \cite{doi:10.1073/pnas.1611835114}, and that rolls the weights back if a held-out retention probe indicates general capability has eroded. Together the three capabilities convert verification from an event that fires once per query into a discipline that compounds across cycles.

A practical consequence of the joint architecture is that the three serving tiers are not static cost categories but a self-organising compute ladder under the cycle loop, and that two distinct statelessness problems of the standard inference paradigm are converted into two stateful mechanisms by two different architectural elements. The first statelessness problem is that retrieval-augmented generation pays the full retrieval-and-generation cost in every session because the model never absorbs what retrieval just taught it; CAEM closes this by feeding verified retrieval-augmented outputs into the self-improvement loop's training pool, so each cycle's low-rank-adapter update folds the previously retrieval-dependent reasoning into the generator's parametric capacity. The next cycle is therefore able to answer a measurable share of the previously retrieval-dependent queries from the zero-shot tier alone, without paying the retrieval cost. The second statelessness problem is that even when the model already knows the answer, the standard protocol re-generates it from scratch on every recurrence and offers no guarantee that two sessions produce the same answer; CAEM closes this through the verified episodic memory, where exact and near-duplicate queries are served directly from the stored answer at embedding-lookup cost, and where the retroactive re-verification pass at every cycle boundary refreshes the stored composite of every cached entry under the locked verifier ensemble and the cycle-boundary-refit composite against the updated generator's outputs, so that the cached answer's quality tracks the model's current capability rather than freezing at the moment of admission. At cycle zero almost every query falls to the retrieval-augmented tier because neither mechanism has had time to accumulate; cycle after cycle, the zero-shot tier absorbs queries that the self-improvement loop has taught the generator to answer, the memory-direct tier absorbs queries that have recurred at least once and whose cached answer still clears the retroactive prune threshold, and the retrieval-augmented tier retains only the genuinely novel queries that neither mechanism has yet covered. The architecture iterates this redistribution until the generator approaches its parametric capacity ceiling, at which point the per-cycle increments enter the bounded stationary band registered in \Cref{ch:methodology}; the six-cycle horizon evaluated in this thesis is a finite, verifiable demonstration of the same pattern on the committed-cycle subsequence, not its asymptote.

The user-facing surface mirrors the same logic. Every served answer carries the calibrated probability of correctness as a continuous percentage alongside the four-class decision tag, so a user sees both the categorical reliability label (verified, provisional, conflicting, or insufficient) and the underlying probability that the answer is right. When the calibrated probability and the grounding evidence both fall short of the registered floors, the system abstains explicitly rather than serving a guessed answer with an inflated confidence label. The architecture is therefore designed against three deployment-time properties at once: hallucination is reduced cycle after cycle through the verifier-gated self-improvement loop, per-query cost falls as queries migrate from the retrieval-augmented tier into the cheaper zero-shot and memory-direct tiers, and the user is never given a confident answer that the system itself could not certify.

There is also a theoretical case for why this combined intervention can succeed even when verification is itself imperfect. The data-purity perspective notes that under reasonable conditions on the verifier balanced accuracy and the base generation accuracy, the purity of the stored pool can exceed the verification accuracy itself, because base-rate dynamics favour the gate whenever the verifier is more likely to admit a correct answer than to admit an incorrect one. The formal statement of this property, together with the conditions under which it implies monotonic improvement across cycles, appears in \Cref{ch:methodology}. For the rationale here, the relevant point is that the architecture does not require a perfect verifier; it requires a verifier whose balanced accuracy exceeds one half on the distribution of claims actually scored, and the empirical chapters report whether that condition is satisfied on the five-benchmark panel.

The scope of the intervention is bounded so that it can be argued empirically within an undergraduate research budget. The thesis evaluates a single small instruction-tuned decoder backbone, namely Qwen-2.5-3B-Instruct, on five factual question-answering benchmarks partitioned into a three-benchmark training panel (FEVER, TriviaQA, CommonsenseQA) and a two-benchmark transfer panel (TruthfulQA, StrategyQA), with multitask language understanding held out as an out-of-distribution retention control alongside test-fold slices of two of the training benchmarks. Three benchmarks registered earlier in the project (HotpotQA~\cite{yang-etal-2018-hotpotqa}, Natural Questions, and ARC-Challenge) were retired before the main run on cycle-zero diagnostic evidence that they violated the verifier-balanced-accuracy precondition for self-improvement; all three are kept on disk as registered exclusion evidence rather than carried forward into the live panel. The Natural Questions training split is also retained inside the retrieval-utility classifier's training distribution to broaden the coverage of retrieval-utility regimes the head must learn; HotpotQA and ARC-Challenge are excluded from every downstream component. The self-improvement trajectory is registered to run for up to ten cycles in line with the equilibrium horizons reported in continual-learning studies, with an equilibrium-saturation gate at every cycle boundary that terminates the trajectory early when the post-rollback subsequence stabilises; the realised trajectory closes at the cycle-five close under that gate. The hardware target is a single consumer-grade graphics processor under sixteen-bit mixed-precision arithmetic, so that the result is reproducible without specialised infrastructure. The next section formalises the problem the intervention is designed to solve, before the objectives, methodology summary, and scope are presented in the remaining sections of this chapter.

\section{Problem Statement}
\label{sec:problem}

\subsection{Formal problem}
\label{sec:formal-problem}

Let a deployed question-answering system be specified by a generator distribution and a retrieval substrate, both held fixed at deployment time. A query arrives and produces an answer through the system's per-query pipeline, after which the system must decide whether to commit the produced answer into a persistent state that will influence future answers. The thesis addresses the following problem. Given a stream of queries drawn from a knowledge-intensive distribution, design a deployable system whose committed answer pool is dominated by correct answers, whose user-facing answers carry a precision guarantee that holds under modest distribution shift between calibration data and live data, whose computational cost per query falls as the system processes more of the stream, and whose general capability does not erode as it accumulates domain-specific knowledge. The system is required to satisfy these properties without per-domain calibration at deployment, without external service calls, and on a single small instruction-tuned generator.

The objects under control are the storage decision applied to each generated answer, the calibration of the confidence used to make that decision, the routing of incoming queries among a memory tier and one or more generation tiers, and the cycle-boundary update of the model parameters. The objects not under control are the generator's prior over factual content, the retrieval corpus, and the eventual stream distribution. Within the controlled set the design must produce a falsifiable empirical claim against a panel of external baselines under matched protocol.

\subsection{The three structural shortcomings the system must address}
\label{sec:three-shortcomings}

The empirical record reviewed in \Cref{sec:background} shows three reinforcing shortcomings of the standard inference paradigm that the present system must address simultaneously.

The first shortcoming is stateless operation without knowledge accumulation. The standard inference protocol processes each query independently. A reasoning chain that produces a correct answer is discarded the moment the answer is delivered, so a substantively similar query reaching the system tomorrow triggers the full reasoning effort again, with no guarantee that it will land on the same answer. The cost is twofold. Wasted computation accumulates against any deployment budget, and inconsistent answers across sessions erode user trust independently of any individual answer's correctness.

The second shortcoming is poor confidence calibration. Modern neural networks routinely assign probability mass to incorrect outputs that exceeds their actual accuracy by a wide margin \cite{pmlr-v70-guo17a, kadavath2022language}, and single-signal correctness predictors operate close to random on the discrimination axis \cite{xiong2023can}. The deployed model therefore cannot reliably tell its user, or its own internal routing logic, when an answer should be trusted. A user-facing system that fails on this axis produces two harmful patterns at once: confidently wrong answers that the user accepts, and uncertainly correct answers that the user rejects. Both patterns hurt deployment value.

The third shortcoming is the absence of self-improvement. Once deployed, large models neither learn from their mistakes nor adapt to the patterns of an institution that uses them. Naive online fine-tuning is not a remedy because it produces catastrophic forgetting on the broad capabilities the base training instilled \cite{doi:10.1073/pnas.1611835114}. The implication is that any verification effort the deployment performs today, whether automated or human, provides no compounding benefit tomorrow. The system that is overconfident in the morning is the same system that is overconfident in the afternoon, with the same failure modes intact.

\subsection{Coupling and what is excluded}
\label{sec:coupling-exclusions}

The three shortcomings reinforce one another. A system without persistent memory cannot retain verified answers, so verification cannot compound. A system without calibrated confidence cannot decide when memory is worth consulting, so even a perfect memory would be under-used. A system without self-improvement cannot translate accumulated verification into better future generation, so calibration alone leaves the underlying error rate unchanged. Addressing any one shortcoming in isolation produces a system that is incrementally better on one axis and unchanged on the other two. The thesis therefore commits to a joint treatment in which memory feeds routing, routing feeds generation, and verification feeds both memory and the cycle-boundary update.

Two classes of approach are explicitly excluded from the present scope. The first is changes to the generator's pre-training corpus or pre-training objective. The thesis treats the generator as fixed and instrument the system around it, so that the architectural claim is separable from any pre-training advantage. The second is reliance on external commercial verification services. Every verification signal the system uses must be reproducible from the open weights of a small specialised model and a publicly available retrieval corpus, so that the result is reproducible without paid infrastructure. The boundary of the problem the thesis solves is therefore an architecture that operates on a fixed open generator, with open verification machinery, on standard hardware, and that satisfies the four properties registered in \Cref{sec:formal-problem} on the five-benchmark factual question-answering panel of \Cref{sec:setup}. The next section converts this problem statement into the specific objectives that organise the rest of the work.

\section{Objective}
\label{sec:objective}

\subsection{Primary objective}
\label{sec:primary-objective}

The primary objective of this thesis is to develop a deployable architecture that reduces confident factual hallucination in knowledge-intensive question answering on a single small instruction-tuned generator, while preserving the generator's general capability across self-improvement cycles. The architecture, named the confidence-aware episodic memory with self-improvement, addresses the three reinforcing shortcomings registered in \Cref{sec:problem} through a verified episodic memory paired with a calibrated probability composite, a four-outcome storage decision, a confidence-aware three-tier router with an independent safety override, retroactive cycle-boundary review of stored content, and a constrained self-improvement loop guarded by a multi-modal retention probe on held-out distributions. The success of the architecture is measured against a panel of external baselines under matched protocol on a five-benchmark factual question-answering panel partitioned into three training benchmarks and two transfer benchmarks, using a composite hallucination metric that aggregates eight measurable subtypes of confident-error failure mode, together with per-tier latency and the per-cycle retention ratio as auxiliary diagnostic axes.

The architecture is a calibration-supervised self-improvement regime: at every cycle boundary a labelled calibration fold refits the per-benchmark calibrated probability composite under a shrinkage prior toward the pooled fit, while the verifier-curated pseudo-labels generated by the per-query pipeline carry the gradient signal at the cycle-boundary fine-tune. The storage gate is a fixed-threshold rule on the calibrated probability that is registered before the main run begins and held constant across cycles; only the composite's per-signal isotonic curves and per-benchmark boost coefficients refit at each cycle boundary, so the operating point of the storage decision is anchored against a single registered scalar that any reader can locate in the configuration. The research trajectory reported in this thesis uses a fixed cold-start fold of 1500 samples that is re-scored under each cycle's post-fine-tune model, while the production deployment specified in \Cref{ch:methodology} constructs a fresh stratified labelled fold from each cycle's accumulated admissions.

\subsection{Specific objectives}
\label{sec:specific-objectives}

The primary objective decomposes into seven specific objectives. Each is a separate engineering or scientific deliverable that contributes to the joint claim and that admits an independent empirical check in the later chapters.

\begin{enumerate}
\item \textbf{Verified episodic memory with a four-outcome decision:} Design an episodic memory whose admission rule is finer than a single binary threshold. A generated answer enters one of four mutually exclusive outcomes determined by a calibrated probability composite, a confabulation early-exit rule that fires when high internal confidence coincides with low external grounding, and a deferral path that holds borderline answers for cycle-boundary reconsideration. Storage uses a similarity-indexed encoder over query and reasoning content with a novelty filter and a value-based pruning policy that combines recency, retrieval frequency, and stored-answer quality.

\item \textbf{Pre-generation confidence, three-tier routing with safety override, and a retrieval-utility classifier:} Compute a pre-generation confidence estimate from a single short forward pass on the query that fuses a short-greedy-decode token-probability aggregate with an encoder-layer convergence ratio, calibrate it via per-benchmark temperature scaling on a disjoint fold, and use it independently from the memory similarity score and the stored-answer quality. The router combines query-to-memory similarity with stored-answer quality into a single dispatch score that selects between memory-direct retrieval, zero-shot generation, and retrieval-augmented generation. An independent safety override forces the conservative tier whenever pre-generation confidence falls below a fixed floor, so that grounded retrieval is available before generation begins and is not entangled with the quality of memory the query happens to land near. A retrieval-utility classifier (Stage 3b) then intercepts the safety-override dispatch and the score-formula fall-through to the retrieval-augmented tier and refines the default into either the zero-shot tier or the retrieval-augmented tier on a per-query basis, fitted under nested cross-validation on a thirty-six-feature query-side and passage-side vector with a registered classifier-side threshold; the empirical contribution of this stage is read off the principal head-to-head ablation reported in \Cref{ch:evaluation}.

\item \textbf{Nine-signal post-generation verifier with a calibrated probability composite:} Combine nine complementary signals into a per-answer record across four families: internal calibration of the generator, sample-set agreement across resampled chains, external grounding against retrieved evidence under three aggregation rules, and question-answer relevance. The signals enter a calibrated probability composite that fits per-signal isotonic regression both on the pooled training-panel fold and on each per-benchmark slice, blends each per-benchmark child curve with the pooled fit through a James-Stein style shrinkage prior, and aggregates the per-signal calibrated probabilities through a regularised logistic regression layer fitted on the same calibration fold. The composite replaces the fixed weighted sums used in earlier verifier ensembles \cite{farquhar2024semantic}.

\item \textbf{Fixed-threshold storage gate with two registered cut-points:} Convert the calibrated probability composite into the storage and deferral decisions through a fixed-threshold rule whose two cut-points and one abstention floor are registered before the main run begins and held constant across cycles. The storage cut-point is selected against the cycle-zero threshold sweep so that the realised per-store empirical precision on the calibration fold is anchored to a known number; the deferral cut-point admits a wider band of marginal entries into the deferred buffer for cycle-boundary reconsideration. The fixed-threshold rule replaces the split-conformal procedure considered earlier in the project, which was rejected on the registered cycle-zero evidence that the calibration-versus-evaluation distribution shift on the panel breaks the marginal-coverage premise of the conformal procedure; the fixed-threshold rule is robust to that shift and surfaces the operating point as a single registered scalar.

\item \textbf{Retroactive review and deferred reconsideration at cycle boundaries:} Because the verifier's internal-calibration signals read the generator's hidden state directly, stored composites are model-relative snapshots that drift after every fine-tune. A retroactive re-verification pass at every cycle boundary rescores stored episodes under the updated model and prunes any entry whose composite has fallen below the storage threshold. A companion pass over the deferred buffer promotes entries that have crossed the storage threshold and discards entries whose age has exceeded a fixed time-to-live.

\item \textbf{Constrained self-improvement with a parameter-efficient adapter and a multi-modal retention guard:} Curate the fine-tuning pool from stored episodes that exceed a strict quality threshold set above the storage cut-point and pass the curated pool through a benchmark-balancing reweighting layer that prevents single-benchmark dominance. Fine-tune the generator through a low-rank adaptation layer on the attention and feed-forward projections rather than through full-parameter fine-tuning \cite{hu2022lora}, with the trainable footprint at roughly two percent of the backbone parameter count and the underlying backbone frozen across the cycle. The bounded adapter parameter budget, layered on a frozen backbone, acts as the implicit cycle-to-cycle drift bound in the spirit of elastic weight consolidation \cite{doi:10.1073/pnas.1611835114}, without the hand-tuned quadratic anchor coefficient that a full-parameter fine-tune would require. Probe a registered set of held-out distributions before and after each fine-tune and abort the cycle with an adapter rollback when the worst probe ratio falls below a fixed floor; the memory is preserved across rollbacks, so accumulated knowledge survives abort events even when the adapter does not.

\item \textbf{Empirical validation against external baselines and architectural ablations:} Evaluate the architecture on a five-benchmark factual question-answering panel partitioned into a three-benchmark training panel (FEVER, TriviaQA, CommonsenseQA) and a two-benchmark transfer panel (TruthfulQA, StrategyQA), with multitask language understanding held out as an out-of-distribution retention control alongside test-fold slices of two of the training benchmarks. Report the composite hallucination metric, the per-tier behaviour, the retention ratio, and the calibration trajectory cycle by cycle. Compare against eight external baselines under matched protocol using paired statistical tests with multiple-comparison correction within each benchmark family, and run pre-registered architectural ablations whose expected effect bands are declared in advance.
\end{enumerate}

\subsection{Pre-registered hypotheses}
\label{sec:hypotheses}

The thesis registers five falsifiable hypotheses ahead of the empirical chapters. Each hypothesis names a specific quantity, a deciding direction, and a deciding statistic, so that the empirical chapters either confirm or refute the hypothesis without ambiguity. The first four hypotheses pair with the load-bearing formal claims registered in \Cref{ch:methodology}; the fifth hypothesis tests the operational claim about deployed cost that is independent of the theoretical scaffolding.

\begin{enumerate}
\item \textbf{Verifier balanced accuracy exceeds one half on the calibration-matched distribution at every cycle:} The hypothesis is the empirical precondition of the pool-purity theorem. The deciding statistic is the verifier balanced accuracy at the fitted storage threshold, measured on the cycle's calibration fold and reported per benchmark. Falsification is any cycle in which the precondition fails on the pooled training panel.

\item \textbf{Pool purity is non-decreasing across cycles whenever the precondition is satisfied:} The hypothesis tests the monotonicity prediction. The deciding statistic is the per-cycle pool purity trajectory, and the deciding direction is non-decreasing in the precondition-satisfied regime. Falsification is any cycle in which the precondition holds and pool purity decreases.

\item \textbf{Late-cycle pool-purity increments enter a bounded stationary band on the committed-cycle subsequence:} The hypothesis tests the convergence prediction in its reframed form. The deciding statistic is the per-cycle pool-purity reading on the committed-cycle subsequence and the realised width of the band across the late portion of the trajectory; the sharper geometric-rate test is registered as deferred future work pending a horizon of at least fifteen committed cycles.

\item \textbf{The composite hallucination metric at the realised cycle horizon approaches a parameter-bounded asymptote:} The hypothesis tests the operational prediction implied by the asymptotic-elimination identity recorded in \Cref{thm:asymptotic-elim}. The deciding statistic is the gap between the measured composite hallucination metric at the final cycle and the predicted asymptote, where the asymptote is bounded by the verifier balanced accuracy and the base accuracy on the relevant fold.

\item \textbf{Per-cycle Tier 1 hit rate grows monotonically as memory accumulates, producing a measurable per-query cost reduction:} The hypothesis tests the operational claim that the architecture reduces deployed cost as the system processes more of the query stream. The deciding statistic is the per-cycle share of incoming queries served by the memory tier and the corresponding per-cycle pooled latency.
\end{enumerate}

The empirical chapters report each hypothesis with its deciding statistic and tag each as supported, partially supported, or unsupported at the trajectory's end. Together the seven specific objectives and the five hypotheses constitute the complete falsifiable content of the architectural claim, each anchored to a measurement reported in \Cref{ch:evaluation}. \Cref{tab:hyp-registration} summarises the registration in compact form.

\begin{table}[!htbp]
\centering
\small
\begin{tabular}{@{}p{0.6cm}p{6.6cm}p{4.5cm}p{2.6cm}@{}}
\toprule
\textbf{ID} & \textbf{Hypothesis} & \textbf{Deciding statistic} & \textbf{Formal claim} \\
\midrule
H1 & Verifier balanced accuracy exceeds one half on the calibration-matched distribution at every cycle. & Per-cycle verifier balanced accuracy at the fitted storage threshold, pooled training panel. & \Cref{thm:purity} precondition \\
H2 & Pool purity is non-decreasing across cycles whenever the precondition is satisfied. & Per-cycle pool-purity trajectory under the precondition-satisfied regime. & \Cref{thm:purity} \\
H3 & Late-cycle pool-purity increments enter a bounded stationary band on the committed-cycle subsequence. & Per-cycle pool-purity reading on the committed-cycle subsequence and realised band width. & \Cref{thm:convergence} \\
H4 & The composite hallucination metric at the realised cycle horizon approaches a parameter-bounded asymptote. & Gap between final-cycle composite hallucination metric and the predicted asymptote. & \Cref{thm:asymptotic-elim} \\
H5 & Per-cycle Tier 1 hit rate grows monotonically as memory accumulates, producing a measurable per-query cost reduction. & Per-cycle Tier 1 share and per-cycle pooled latency. & Operational claim (\Cref{sec:econ-deployment}) \\
\bottomrule
\end{tabular}
\caption[The five pre-registered hypotheses with deciding statistic and the theoretical or operational claim each tests.]{The five pre-registered hypotheses with deciding statistic and the theoretical or operational claim each tests. Verdicts are reported in \Cref{tab:hypothesis-verdicts}.}
\label{tab:hyp-registration}
\end{table}

\section{Methodology in Brief}
\label{sec:methodology-brief}

\subsection{Architecture overview}
\label{sec:arch-overview}

The architecture organises a single query into an eight-stage pipeline that is traversed once per query and wrapped in a per-cycle update loop. The first stage produces a pre-routing confidence from a single short forward pass over the query, fusing a short-greedy-decode token-probability aggregate with an encoder-layer convergence ratio. The second stage encodes the query with a sentence-embedding encoder and searches a similarity-indexed episodic memory for the nearest stored entry, returning a similarity score and the neighbour's stored-quality score. The third stage combines the similarity score and the stored-quality score into a single dispatch score and routes the query to one of three tiers, with an independent safety override that forces the conservative tier whenever the pre-routing confidence falls below a fixed floor. A retrieval-utility classifier then intercepts both the safety-override dispatch and the score-formula fall-through to the retrieval-augmented tier and refines that default into either the zero-shot tier or the retrieval-augmented tier on a per-query basis, so that retrieval is consulted only on queries it is expected to help; the classifier is specified in \Cref{sec:ruc}. The fourth stage executes the chosen tier: the memory-direct tier serves the neighbour's stored answer; the zero-shot tier runs a generation on the fine-tuned generator; the retrieval-augmented tier retrieves passages from a public corpus and runs a grounded generation. The fifth stage passes every generated answer through a nine-signal post-generation verifier whose four signal families (internal calibration of the generator, sample-set agreement across resampled chains, external grounding under three aggregation rules, and question-answer relevance) feed a calibrated probability composite that converts each signal into a calibrated probability via per-signal isotonic regression and aggregates them through a regularised logistic regression layer fitted on a disjoint calibration fold. The sixth stage applies a four-outcome decision (store, defer, abstain, discard) to the composite output, with a confabulation early-exit rule that fires when high internal confidence coincides with low external grounding and takes precedence over the composite thresholds. The seventh stage executes the implied memory and user-facing actions. The eighth stage runs at cycle boundaries and is detailed below. \Cref{fig:architecture} gives the high-level view; the full per-stage pipeline appears in \Cref{ch:methodology}.

\begin{figure}[!htbp]
\centering
\resizebox{\textwidth}{!}{%
\begin{tikzpicture}[
    >=stealth,
    node distance=8mm and 10mm,
    box/.style={rectangle, rounded corners=4pt, draw=#1!70!black, fill=#1!12,
                line width=0.9pt, align=center, text width=3.0cm, minimum height=1.6cm,
                inner sep=4pt, font=\normalsize\bfseries},
    store/.style={rectangle, rounded corners=4pt, draw=teal!70!black, fill=teal!10,
                  line width=0.9pt, align=center, text width=3.0cm, minimum height=1.6cm,
                  inner sep=4pt, font=\normalsize\bfseries},
    arrow/.style={->, line width=1.1pt, draw=#1!70!black},
    darrow/.style={->, line width=1.0pt, draw=#1!70!black, dashed},
    elabel/.style={font=\small, fill=white, inner sep=2pt, text=black!75}
]
%% Top row
\node[box=blue]                                 (qry)  {Query};
\node[box=cyan,   right=of qry]                 (pre)  {Stage 1\\Pre-route};
\node[box=teal,   right=of pre]                 (enc)  {Stage 2\\Encode};
\node[box=orange, right=of enc]                 (route){Stage 3\\Route};

%% Memory above encode
\node[store, above=of enc]                      (mem)  {Memory};

%% RUC (Stage 3b) between Route and Tier band
\node[box=yellow!60!orange, right=8mm of route, fill=yellow!18, text width=2.4cm, minimum height=1.4cm]
                                                (ruc)  {Stage 3b\\RUC};

%% Tier band right of RUC
\node[box=green,  right=10mm of ruc, yshift=18mm]  (t1) {Tier 1\\memory};
\node[box=yellow, right=10mm of ruc]                (t2) {Tier 2\\zero-shot};
\node[box=red,    right=10mm of ruc, yshift=-18mm] (t3) {Tier 3\\RAG};

%% External corpus
\node[box=blue, dashed, below=8mm of t3, fill=blue!4] (kb) {Public\\corpus};

%% Verify, Decide, Output
\node[box=violet, right=22mm of t2]   (ver) {Stage 5\\Verify};
\node[box=purple, below=of ver]       (dec) {Stage 6\\Decide};
\node[box=green,  right=of ver]       (out) {Stage 7\\Output};

%% Generator below route
\node[store, below=12mm of route]     (gen) {Generator};

%% Arrows -- ingress
\draw[arrow=blue]   (qry)   -- (pre);
\draw[arrow=cyan]   (pre)   -- (enc);
\draw[arrow=teal]   (enc)   -- (route);

%% Memory probe
\draw[darrow=teal]  (enc) -- (mem);
\draw[darrow=teal]  (mem) -| (route);

%% Router to T1 direct; T2/T3 fall-through via RUC
\draw[arrow=green, rounded corners=4pt]  (route.east) -- ++(0.3,0) |- (t1.west);
\draw[arrow=orange] (route) -- (ruc);
\draw[arrow=orange] (ruc.east) to[out=20, in=180]  (t2.west);
\draw[arrow=red]    (ruc.east) to[out=-20, in=180] (t3.west);

%% Safety override (pre-route low confidence -> RUC, which then decides T2 vs T3)
\draw[darrow=red, rounded corners=4pt]   (pre.south) -- ++(0,-0.8) -| node[elabel, pos=0.25, anchor=north] {safety} (ruc.south);

%% Memory to T1, Generator to T2
\draw[arrow=green]  (mem.east) to[out=0, in=120] (t1.north);
\draw[arrow=orange] (gen.east) to[out=0, in=-120] (t2.south);

%% External corpus to T3
\draw[darrow=blue]  (kb) -- (t3);

%% Tiers to verify
\draw[arrow=green]  (t1) -| (ver.north);
\draw[arrow=orange] (t2) -- (ver);
\draw[arrow=red]    (t3) -| (ver.south);

%% Verify -> Decide -> Output
\draw[arrow=violet] (ver) -- (dec);
\draw[arrow=green]  (ver) -- (out);

%% Decide -> Memory (commit), Memory -> Generator (training pool)
\coordinate (commitUp) at (0, 4.2);
\draw[darrow=purple, rounded corners=4pt] (dec.east) -- ++(4.5, 0) |- node[elabel, pos=0.25] {commit} (mem.north |- commitUp) -- (mem.north);
\draw[darrow=purple, rounded corners=4pt] (mem.south east) -- ++(0.4,0) |- node[elabel, pos=0.75] {train} (gen.west);

%% Stage 8 frame
\node[draw=purple!65, dashed, line width=0.7pt, rounded corners=4pt,
      fit={(mem) (gen) (route)}, inner sep=4pt,
      label={[font=\scriptsize\bfseries, text=purple!75!black]above:Stage 8 \ Self-improvement (cycle boundary)}] {};
\end{tikzpicture}
}

\caption[Architecture overview]{\small Eight-stage pipeline: pre-route, encode, route, Stage 3b retrieval-utility classifier on the score-formula and safety-override fall-through to the retrieval-augmented tier, tier execution, verify, decide, output, and a cycle-boundary self-improvement block; full specification in \Cref{ch:methodology}.}
\label{fig:architecture}
\end{figure}

\subsection{Per-query pipeline}
\label{sec:per-query-pipeline}

A query enters the system and is encoded into a sentence-embedding vector. The encoded vector probes the episodic memory and returns the cosine similarity to the nearest stored entry together with the stored-quality score of that neighbour. In parallel, the pre-routing confidence is computed from a single short forward pass on the query that fuses a short-greedy-decode token-probability aggregate with an encoder-layer convergence ratio, and is calibrated through per-benchmark temperature scaling on a disjoint fold. The router combines the similarity and stored-quality scores into a single dispatch score and selects the memory-direct tier when the score is high, the zero-shot tier in a middle band, and the retrieval-augmented tier when the score is low. An independent safety override forces the retrieval-augmented tier whenever the pre-routing confidence falls below a fixed floor, so that low pre-generation readiness produces grounded retrieval before any generation happens. A retrieval-utility classifier then intercepts the safety-override dispatch and the score-formula fall-through to the retrieval-augmented tier and refines the default to either the zero-shot tier or the retrieval-augmented tier on a per-query basis, using a thirty-six-feature query-side and passage-side feature vector and a registered classifier-side threshold; on queries the classifier judges retrieval-hurt the dispatch is redirected to the zero-shot tier, on queries the classifier judges retrieval-helpful the retrieval-augmented dispatch stands, with the full specification in \Cref{sec:ruc}. The chosen tier produces an answer, which then enters the verifier.

The verifier evaluates nine complementary signals across four families. Internal calibration signals read the generator's own state and include token-level uncertainty, dropout-based disagreement, and an internal-state composite. Sample-set signals measure agreement across resampled chains, including average pairwise similarity, chain-to-answer entailment, and normalised semantic entropy \cite{farquhar2024semantic}. External grounding signals measure entailment between retrieved passages and the answer under three aggregation rules: top-one entailment, pooled mean entailment, and a weakest-link atomic-fact decomposition that fires only when the answer is long enough to admit decomposition. A question-answer relevance signal closes the off-topic-but-grounded failure mode by checking whether the answer addresses the question. The signals enter a calibrated probability composite that maps each per-signal score to a calibrated probability via isotonic regression and aggregates them through a regularised logistic regression layer fitted on a disjoint calibration fold; the calibrated output is the storage score. A four-outcome decision then maps the storage score, together with a confabulation early-exit rule, to one of four mutually exclusive actions: commit to memory, defer to a bounded buffer with a time-to-live, abstain explicitly, or discard. The user-facing action that follows depends on the outcome, so that the user receives a stored answer, a generated answer, or a calibrated refusal rather than a silently uncertain response.

\subsection{Calibration discipline}
\label{sec:calibration-discipline}

The storage score is converted into the storage and deferral decisions through a fixed-threshold rule whose two cut-points are registered before the main run begins and held constant across cycles. The storage cut-point is selected against the cycle-zero threshold sweep so that the realised per-store empirical precision on the calibration fold is anchored to a known number; the deferral cut-point admits a wider band of marginal entries into the deferred buffer. The fixed-threshold rule replaces an earlier split-conformal formulation that fitted both thresholds at pre-registered conformal levels and re-fitted them at each cycle boundary under exponential-moving-average smoothing; that formulation was rejected on the registered cycle-zero evidence that the calibration-versus-evaluation distribution shift on the five-benchmark panel breaks the marginal-coverage premise of the conformal procedure, with the realised eval-fold precision falling fifteen to fifty percentage points below the conformal target on training benchmarks. The architectural cut-points are therefore frozen across the trajectory; only the calibrated probability composite refits at each cycle boundary on the cycle's own labelled fold, so the operating point of the storage decision does not drift as the underlying score distribution shifts.

\subsection{Cycle-boundary self-improvement loop}
\label{sec:cycle-loop}

The cycle loop iterates the per-query pipeline over a fixed workload, accumulates verified episodes into memory, and runs three cycle-boundary passes before fine-tuning. The retroactive re-verification pass rescores every stored episode under the updated generator, promotes entries whose composite has crossed the storage threshold, and prunes entries whose composite has fallen below it. The deferred reconsideration pass examines every entry in the bounded deferral buffer and either promotes, expires, or requeues the entry. The memory consolidation pass collapses near-duplicate stored episodes that exceed a similarity threshold, picking the highest-quality representative within each near-duplicate cluster. Fine-tuning then draws a training pool from stored episodes whose composite score exceeds a strict quality threshold above the storage threshold, ensuring the training pool is held to a higher bar than the user-facing retrieval tier. The fine-tune is regularised implicitly through the bounded low-rank-adapter parameter budget layered on a frozen backbone, in the spirit of elastic weight consolidation \cite{doi:10.1073/pnas.1611835114} but without an explicit quadratic anchor coefficient. Before and after the fine-tune, the cycle probes a held-out distribution and computes a retention ratio. When the retention ratio falls below a fixed floor, the cycle aborts: weights are rolled back to the previous checkpoint and the retroactive pass for the next cycle is skipped. The memory is never rolled back, so accumulated knowledge survives abort events even when weights do not. This asymmetry decouples weight-level drift from memory-level accumulation and is what makes the self-improvement loop safe to iterate.

\Cref{fig:experimental-workflow} illustrates the trajectory across early, mid, and late phases. Each phase runs the same per-query pipeline over a fresh workload of queries, then applies the cycle-boundary update passes under the retention guard. The three boxes are representative phases rather than literal per-cycle depictions; the actual trajectory is reported cycle by cycle in \Cref{ch:evaluation}.

\begin{figure}[H]
\centering
\begin{tikzpicture}[
    >=Stealth,
    font=\sffamily,
    %% phase card -- left column, with chip header and bullet list
    phase/.style 2 args={rectangle, rounded corners=6pt,
        draw=#1!75!black, line width=0.9pt, fill=#1!10,
        align=left, text width=6.6cm, minimum height=2.8cm,
        inner sep=7pt, font=\small,
        blur shadow={shadow blur steps=4, shadow blur extra rounding=1pt,
                     shadow xshift=0.5ex, shadow yshift=-0.5ex, shadow opacity=20}},
    %% cycle-boundary band -- thin grey strip between phases
    boundary/.style={rectangle, rounded corners=3pt,
        draw=gray!50, line width=0.7pt, fill=gray!8,
        align=center, text width=6.6cm, minimum height=0.95cm,
        inner sep=4pt, font=\footnotesize\itshape, text=black!75},
    %% logs callout -- right column, mirrors phase card colour
    logs/.style 2 args={rectangle, rounded corners=4pt,
        draw=#1!55!black, line width=0.7pt, fill=#1!6,
        align=left, text width=4.8cm, minimum height=2.8cm,
        inner sep=5pt, font=\scriptsize},
    %% arrow between phases through boundary band
    flow/.style={->, line width=1.4pt, draw=#1!70!black}
]

%% ----- Phase cards (left column) ----------------------------------------
\node[phase={red}{}] (c1) {%
\colorbox{red!75!black}{\color{white}\strut\,\textbf{\scriptsize EARLY CYCLES}\,}\quad\textit{\footnotesize cold start}\\[4pt]
\(\bullet\) Memory: \textbf{empty $\rightarrow$ sparse}\\
\(\bullet\) Most queries to Tier 2 / Tier 3\\
\(\bullet\) Nine-signal verify $\rightarrow$ four-outcome decide\\
\(\bullet\) Store / Defer / Abstain / Discard%
};

\node[boundary, below=0.32cm of c1] (cons1) {%
Cycle boundary: retroverify \(+\) defer-reconsider \(+\) consolidate \\
low-rank-adapter fine-tune \(+\) retention guard
};

\node[phase={orange}{}, below=0.32cm of cons1] (c2) {%
\colorbox{orange!75!black}{\color{white}\strut\,\textbf{\scriptsize MID CYCLES}\,}\quad\textit{\footnotesize accumulating}\\[4pt]
\(\bullet\) Memory: \textbf{growing}\\
\(\bullet\) Tier 1 share rises\\
\(\bullet\) Training pool above storage threshold\\
\(\bullet\) Retroverify promotes/prunes under updated model%
};

\node[boundary, below=0.32cm of c2] (cons2) {%
Cycle boundary: same passes; on retention failure, weights roll back, memory persists
};

\node[phase={green}{}, below=0.32cm of cons2] (c3) {%
\colorbox{green!75!black}{\color{white}\strut\,\textbf{\scriptsize LATE CYCLES}\,}\quad\textit{\footnotesize near equilibrium}\\[4pt]
\(\bullet\) Memory: \textbf{mature, near-equilibrium}\\
\(\bullet\) Highest Tier 1 share\\
\(\bullet\) Composite evaluation score reported\\
\(\bullet\) Composite hallucination metric vs.\ baselines%
};

%% ----- Logs callouts (right column) -------------------------------------
\node[logs={red}{}, right=1.0cm of c1] (m1) {%
\textbf{Per-cycle logs}\\[2pt]
\(\bullet\) (store, defer, abstain, discard) counts\\
\(\bullet\) training-pool size\\
\(\bullet\) routing distribution by tier%
};

\node[logs={orange}{}, right=1.0cm of c2] (m2) {%
\textbf{Cycle-boundary logs}\\[2pt]
\(\bullet\) retroverify (promotions, prunes)\\
\(\bullet\) deferred (promote, drop, keep)\\
\(\bullet\) consolidation clusters resolved\\
\(\bullet\) retention ratio against held-out probe%
};

\node[logs={green}{}, right=1.0cm of c3] (m3) {%
\textbf{Final report \textnormal{(\Cref{ch:evaluation})}}\\[2pt]
\(\bullet\) composite evaluation score, geometric mean\\
\(\bullet\) composite hallucination metric reduction\\
\(\bullet\) per-cycle trajectory \(+\) classifier on/off ablation%
};

%% ----- Vertical flow arrows --------------------------------------------
\draw[flow=red]    (c1)    -- (cons1);
\draw[flow=orange] (cons1) -- (c2);
\draw[flow=orange] (c2)    -- (cons2);
\draw[flow=green]  (cons2) -- (c3);

\end{tikzpicture}
\caption[Experimental workflow]{Self-improvement trajectory across early, mid, and late phases of the realised six-cycle horizon. Each phase runs the per-query pipeline and, at cycle boundaries, the retroactive re-verification, deferred reconsideration, consolidation, and adapter fine-tune under the multi-modal retention guard. Full specification appears in \Cref{ch:methodology}; headline results in \Cref{ch:evaluation}.}
\label{fig:experimental-workflow}
\end{figure}

\section{Scopes and Challenges}

\subsection{Research scope}
\label{sec:scope-research}

The architecture targets factual question answering, defined as the task of producing answers that admit a definite correctness verdict against a reference. The five evaluation benchmarks adopted in this work span factual claim verification (FEVER), open-domain trivia (TriviaQA), multiple-choice commonsense reasoning (CommonsenseQA), truthfulness on common-misconception probes (TruthfulQA), and multi-step commonsense reasoning (StrategyQA), partitioned into a three-benchmark training panel and a two-benchmark transfer panel. Each benchmark admits an objective scoring rule, namely exact match or label accuracy, depending on the task family. A multi-modal retention probe panel is reserved as held-out retention control rather than as a target benchmark: the fine-tuning loop is not optimised on any probe, and the retention guard aborts any cycle in which the worst per-probe accuracy ratio falls below a fixed fraction of the cycle-zero baseline. The probe panel comprises a general-knowledge multiple-choice surface (MMLU), an open-text factoid surface drawn from the TriviaQA test fold, and a commonsense multiple-choice surface drawn from the CommonsenseQA test fold; the worst-probe rule is the registered abort criterion.

The scope explicitly excludes tasks whose evaluation is intrinsically subjective, including creative generation, open-ended dialogue, and stylistic rewriting. Tasks that require knowledge beyond the pre-training cutoff, such as breaking news, market data, and recently issued documents, are excluded because the system does not include a continuous-update mechanism for ingesting external knowledge sources. Code generation and formal mathematical reasoning are excluded so that the evaluation surface remains within natural-language reasoning and the verifier signals retain a uniform interpretation across benchmarks. The system operates on English text only; multilingual transfer and code-switched evaluation are deferred to future work.

Resource scope is set by the engineering envelope of an undergraduate thesis. The generator is a three-billion-parameter open-weight instruction-tuned decoder-only model, chosen for two reasons. First, it admits parameter-efficient fine-tuning on a single consumer-grade graphics processor with thirty-two gigabytes of memory through a low-rank adaptation layer on the attention and feed-forward projections, with the trainable footprint at roughly two percent of the backbone parameter count and the underlying backbone frozen across the cycle; full-parameter fine-tuning on the same backbone size does not fit the engineering envelope and is registered as the rejected alternative on the joint retention-and-feasibility axis. Second, its baseline accuracy on the five benchmarks of the registered panel is high enough that the verifier's expected false-positive rate on stored episodes can be bounded below the level at which a self-improvement loop turns counterproductive. The ten-cycle horizon is sized to capture the early-cycle calibration phase, the mid-cycle accumulation phase, and the late-cycle equilibrium phase predicted by the convergence analysis in \Cref{ch:methodology}, while remaining feasible within a twelve-month compute envelope. The episodic memory is sized to a target capacity in the millions of entries with a tiered index that promotes from a flat structure to a quantised inverted-file structure above a threshold; the workloads reported in this thesis grow the memory to the low tens of thousands of entries, well within the capacity envelope and sufficient to exercise the novelty filter and the value-score pruning rule without saturating them.

\subsection{Technical challenges}
\label{sec:scope-challenges}

The implementation faces six primary technical challenges that the architecture must address.

\begin{itemize}
    \item \textbf{Verifier reliability and stored-quality drift:} The nine-signal post-generation verifier produces a calibrated storage score that doubles as a retrieval-time quality proxy and as a quality filter on the training pool. Three of the nine signals belong to the internal-calibration family and read the generator's hidden state directly, so their distributions shift after every fine-tune. The stored composite of every episode is therefore a model-relative snapshot, not a permanent attribute of the entry. The retroactive re-verification pass at each cycle boundary refreshes every stored composite under the updated generator and applies an upward-only update rule that promotes entries whose composite has crossed the storage threshold and prunes those that have fallen below it. The cost of this discipline is a full re-scoring of every stored episode at each cycle boundary, an order of magnitude smaller than fine-tuning itself but not negligible at scale.

    \item \textbf{Catastrophic forgetting under iterative fine-tuning:} Fine-tuning on domain-specific verified examples risks degrading the generator's general-purpose capabilities. The architecture mitigates this risk through three layers of defence: a strict training-pool quality filter that holds the training pool to a higher bar than the user-facing storage tier, a bounded low-rank-adapter parameter envelope layered on a frozen backbone that acts as the implicit cycle-to-cycle drift bound (in the spirit of elastic weight consolidation \cite{doi:10.1073/pnas.1611835114} but without an explicit quadratic anchor coefficient on the shipped path), and a multi-modal retention probe panel whose worst per-probe ratio gates the commit. The retention probe panel spans a general-knowledge multiple-choice surface, an open-text factoid surface drawn from a training-benchmark test fold, and a commonsense multiple-choice surface drawn from a second training-benchmark test fold; the cycle aborts when any per-probe ratio falls below the fixed floor, weights are rolled back to the previous checkpoint, and the memory is preserved so that a fine-tune that fails in any given cycle does not lose the verified episodes accumulated during it. The empirical contribution of each layer is verified through the architectural ablations reported in \Cref{ch:evaluation}, of which the retrieval-utility-classifier on/off contrast on the realised six-cycle trajectory is the principal landed comparison.

    \item \textbf{Pre-routing confidence calibration:} The pre-routing confidence is used both as a one-dimensional input to the dispatch score and as the sole input to the safety override that forces grounded retrieval, so its calibration influences the system's behaviour twice. The architecture fuses two cheap signals from a single short forward pass, namely a short-greedy-decode token-probability aggregate that scores the first thirty-two tokens of a draft answer and an encoder-layer convergence ratio computed from the forward pass over the query, and temperature-scales the fused scalar per benchmark on a disjoint fold by minimising binary negative-log-likelihood. The non-trivial design choice is that the pre-routing signal does not require the full multi-chain sampling that the post-generation verifier uses, so its cost is paid per query regardless of which tier eventually serves the answer; the cost of the more expensive sample-based signals is deferred to the post-generation verifier on the zero-shot and retrieval-augmented paths only.

    \item \textbf{Memory retrieval quality and novelty filtering:} The episodic memory uses approximate nearest-neighbour search, which raises the possibility of missing a semantically similar entry whose embedding happens to fall near a partition boundary. Near-duplicate writes are filtered at storage time by a novelty threshold on cosine similarity, and the consolidation pass at each cycle boundary collapses near-duplicate clusters by retaining only the highest-quality representative within each cluster. The novelty threshold is a tuning knob with a clear failure mode in either direction: a threshold set too high admits paraphrase-level duplicates that bloat the index without adding coverage, while a threshold set too low rejects legitimate paraphrases that would have improved retrieval. The chosen value is fixed before the main run and held constant across cycles, with sensitivity analyses reported in \Cref{ch:evaluation}.

    \item \textbf{Baseline performance dependency for self-improvement:} The self-improvement argument depends on the base generator performing meaningfully above random on the target benchmarks, because training on verified-but-imperfect episodes can only compound accuracy when the verifier's expected false-positive rate is low enough that the training pool is cleaner than the generator's untrimmed outputs. The risk is mitigated by selecting an instruction-tuned backbone with competitive baseline accuracy on factual question answering, by selecting benchmarks on which the base generator clears that threshold, and by composing the storage decision from the calibrated probability composite, the confabulation early-exit gate, and the in-distribution and out-of-distribution training-pool gate. Each component contracts the storage class below the rate that an unrestricted score would admit, and the conjunction is what makes the precision floor on the storage tier achievable.

    \item \textbf{Evaluation metric selection:} Measuring hallucination reduction requires a metric design that is more discriminative than aggregate accuracy. The composite hallucination metric used in this thesis is the equal-weighted mean over an eight-axis failure-mode taxonomy that includes confident confabulation, factual fabrication, logical fabrication, off-topic answers, defensive evasion, template leakage, false refusal, and length-padded over-generation. A ninth axis is retained in the taxonomy but excluded from the default denominator because the corresponding subtype is structurally unreachable under the verifier backend used here, and its rate is reported separately for completeness. A uniform reduction across eight structurally distinct failure modes is a stronger architectural claim than a larger reduction on any single axis. The composite evaluation score is a geometric mean across accuracy, epistemic quality, retention, calibration, and verifier reliability; the geometric form penalises any axis that collapses, which matches the design constraint that no single axis may be sacrificed for the others. Comparison against the baseline panel is conducted under matched protocol, with identical samples, prompts, and scoring rules, so that effects are not confounded by re-implementation choices.
\end{itemize}

\subsection{Mitigation strategies}
\label{sec:scope-mitigations}

Each challenge above admits a structural mitigation that is built into the architecture rather than added as a downstream patch. Verifier reliability and stored-quality drift are absorbed by the retroactive re-verification pass with its upward-only update rule. Catastrophic forgetting is bounded by the conjunction of the strict training-pool quality filter, the bounded low-rank-adapter parameter envelope on a frozen backbone, and the multi-modal retention guard with adapter rollback. Pre-routing confidence calibration is handled by fusing a short-greedy-decode token-probability aggregate with an encoder-layer convergence ratio computed from a single forward pass, followed by per-benchmark temperature scaling on a disjoint fold. Memory retrieval quality and novelty are managed by the novelty threshold at write time and by the value-score pruning rule that triggers only when capacity is exceeded. Baseline-performance dependency is managed by backbone and benchmark selection together with the confabulation early-exit gate and the in-distribution training-pool gate that excludes transfer-benchmark episodes from the gradient signal, both of which act as additional contraction rules on the storage class. Evaluation-metric risk is managed by reporting both the composite evaluation score (a geometric mean across accuracy, epistemic quality, retention, calibration, and verifier reliability, which penalises any collapsing axis by multiplicative composition rather than masking it under an additive average) and the composite hallucination metric across the eight-axis taxonomy, so that no single collapsing axis can mask a regression on another. The empirical receipts for each mitigation are reported in \Cref{ch:evaluation}.

Although the problems are significant, they can be managed within the resources and the time frame. The aim is to improve the system clearly and measurably while making the architectural commitments that the current state of the literature supports.

\section{Chapter signpost}
\label{sec:ch1-signpost}

This chapter has registered the deployment-safety problem, the architectural opportunity that closes the stateless-overconfident-static triad, the seven specific objectives that decompose the primary objective, and the five pre-registered hypotheses that the empirical chapters will adjudicate. \Cref{ch:literature-review} establishes the technical foundations on which the architectural commitments rest (the transformer backbone, the hallucination taxonomy, natural language inference, sentence embeddings, self-consistency decoding, and catastrophic forgetting) and surveys the empirical literature across hallucination measurement, verification, confidence estimation, memory-augmented architectures, continual learning, and evaluation benchmarks. The survey locates the architectural contribution of this thesis against the prior work that the seven specific objectives draw from, so the methodology that follows in \Cref{ch:methodology} can be read as a single integrated response to a literature whose individual contributions do not close the joint gap.